\section{Innovation \& Solution Design}\label{sec:solution-design}

\subsection{Proposed Innovation: Fair Trade AI with Provenance Tracking}\label{sec:fair-trade-ai}
To address the dilemma between generative AI innovation, intellectual property protection, and the spread of deepfakes, our team proposes an innovative yet realistic solution: a \textbf{Fair Trade AI framework} supported by \textbf{digital provenance and watermarking mechanisms}. Rather than relying on either a total ban or unrestricted market-driven development, this framework combines technical safeguards with governance measures to make generative AI more accountable.

The proposed framework has three main components:
\begin{itemize}
    \item \textbf{Licensed and Opt-in Training Datasets:} AI developers should move toward licensed and opt-in datasets, especially for commercial systems. Instead of indiscriminate web scraping, developers would use datasets obtained through explicit permission, contractual licensing, or lawful data partnerships. This allows creators to decide whether their works can be used for model training and under what conditions.

    \item \textbf{Auditable Licensing and Royalty System:} When creators choose to participate, their works can be registered through a transparent tracking mechanism that records usage rights and supports compensation models. This system does not need to rely exclusively on blockchain; what matters is that it is verifiable, transparent, and capable of building trust between AI developers and content creators.

    \item \textbf{Robust Provenance and Watermarking Mechanisms:} AI-generated content should include provenance signals, such as watermarking, metadata standards, or content authenticity mechanisms. These tools would not make detection perfect, but they would make it easier for platforms, journalists, and users to identify when content is AI-generated. In practice, this would improve traceability and reduce the ease with which synthetic media can be passed off as authentic human-created content.
\end{itemize}

\subsection{Justification and Improvements}\label{sec:justification}
This proposal improves on conventional responses in several important ways.

First, compared with a \textbf{total ban}, the Fair Trade AI framework is more realistic because it allows beneficial innovation to continue. Generative AI still has legitimate applications in education, accessibility, design assistance, and business productivity. A blanket prohibition would likely suppress these benefits along with the harmful uses.

Second, compared with an \textbf{unregulated innovation-first approach}, this framework introduces accountability at the design stage rather than treating harm as an acceptable side effect of rapid deployment. By requiring lawful data sourcing, clearer usage rights, and provenance mechanisms, it reduces the legal and ethical risks that currently make many organizations hesitant to adopt generative AI in professional settings.

Third, compared with relying only on \textbf{post-publication moderation}, this solution is more preventive. It does not wait until harmful content has already spread. Instead, it builds safeguards into the data pipeline, model governance process, and output verification stage. This makes misuse easier to detect, discourage, or limit earlier.

Overall, the proposed framework is more sustainable because it preserves the benefits of generative AI while reducing foreseeable harm to creators, individuals, and the broader public.

\subsection{Comparative Evaluation of Alternative Approaches}\label{sec:comparison}
To clarify why the proposed Fair Trade AI framework is more balanced than conventional responses, \cref{tab:solution-comparison} compares the four main approaches introduced in \cref{sec:alternatives} across four criteria: innovation potential, legal safety, stakeholder protection, and practical feasibility.

\begin{table}[H]
\centering
\caption{Comparison of Alternative Approaches to Generative AI Governance}\label{tab:solution-comparison}
\small
\begin{tabularx}{\textwidth}{@{}P{3cm}P{2.8cm}P{2.5cm}P{2.8cm}L@{}}
\toprule
\textbf{Option} & \textbf{Innovation Potential} & \textbf{Legal Safety} & \textbf{Stakeholder Protection} & \textbf{Practical Feasibility} \\
\midrule
\textbf{Total ban on commercial generative AI} & Low -- blocks both harmful and beneficial uses & High in principle & High for creators and identity protection, but limits broader user benefits & Low -- difficult to enforce consistently across markets and jurisdictions \\
\addlinespace
\textbf{Innovation-first approach with minimal restrictions} & High -- maximizes speed and experimentation & Low -- high exposure to copyright, privacy, and misuse risks & Low -- places creators, individuals, and the public at greater risk & High in the short term, but weak in long-term sustainability due to legal and social backlash \\
\addlinespace
\textbf{Post-publication platform moderation} & Medium to High -- allows continued development & Medium -- some harmful outputs can be removed after release & Medium -- offers partial protection, but often only after damage has occurred & Medium -- more realistic than a total ban, but fundamentally reactive \\
\addlinespace
\textbf{Fair Trade AI with provenance tracking and watermarking} & High -- preserves innovation while introducing accountability & High -- supports lawful data use and clearer responsibility & High -- improves protection for creators, individuals, and the public & Medium -- requires investment, coordination, and compliance effort, but remains realistic for regulated deployment \\
\bottomrule
\end{tabularx}
\end{table}

As shown in \cref{tab:solution-comparison}, the proposed Fair Trade AI framework is not the cheapest or fastest option, but it offers the most balanced combination of innovation, legal defensibility, and stakeholder protection. This makes it a more sustainable and professionally responsible solution than either a complete ban, unrestricted development, or purely reactive moderation.

\subsection{Accepted Trade-offs}\label{sec:trade-offs}
Implementing this framework requires several trade-offs.

\begin{itemize}
    \item \textbf{Higher Compliance and Operational Cost:} Building licensed datasets, negotiating usage rights, maintaining audit systems, and implementing provenance mechanisms all require additional investment. As a result, product development may become slower and more expensive than under a purely growth-driven model.

    \item \textbf{Reduced Dataset Scale and Possible Limits in Diversity:} If developers rely more heavily on opt-in or licensed data, the amount of training material may initially be smaller than what is available through unrestricted web scraping. This may affect model performance, reduce coverage of certain styles or languages, and create new risks of representational imbalance.

    \item \textbf{Greater Friction in User Experience and Deployment:} Verification layers, disclosure requirements, and provenance checks may introduce latency or reduce convenience in some real-time applications. Although these burdens may be unpopular in the short term, they represent a deliberate sacrifice of speed in exchange for stronger trust and accountability.
\end{itemize}

\subsection{Residual Risks and Limitations}\label{sec:residual-risks}
Even with these safeguards, the solution does not eliminate all risks.

\begin{itemize}
    \item \textbf{Watermarking and Provenance Tools Are Not Foolproof:} Determined actors may attempt to strip, degrade, or bypass these protections through model modification, re-encoding, or post-processing techniques. Therefore, provenance should be treated as a risk-reduction tool rather than a perfect guarantee.

    \item \textbf{Open-source or Unofficial Model Redistribution:} Even if leading companies adopt responsible standards, malicious users may fine-tune or adapt open models for harmful purposes outside regulated environments. This means governance cannot rely only on responsible firms; platform policies and legal enforcement still matter.

    \item \textbf{The ``Analog Hole'' Vulnerability:} Someone may capture AI-generated media from a screen using an external recording device, weakening or removing machine-readable provenance signals while preserving deceptive content for human viewers. This shows that technical controls alone cannot fully solve the problem.

    \item \textbf{Dependence on Industry Coordination and Regulatory Support:} A responsible framework will be harder to implement if only a few companies follow it while competitors continue using cheaper and less accountable methods. For this reason, the solution is most effective when supported by shared standards, platform enforcement, and legal incentives.
\end{itemize}

\subsection{Conclusion}\label{sec:solution-conclusion}
Overall, the Fair Trade AI framework offers a more balanced solution than either banning generative AI entirely or allowing it to develop without meaningful constraints. Its main strength is that it combines innovation with responsibility: it preserves the productive and creative benefits of generative AI while introducing more credible protections for creators, individuals, and public trust.

Although the framework involves higher costs, slower deployment, and imperfect technical safeguards, these trade-offs are acceptable because they address the core risks identified in the earlier dilemma analysis (\cref{sec:critical-analysis}). Therefore, this proposal represents a realistic and professionally responsible direction for the development of generative AI systems.
